

\illustration{0.4}{
  Garique305c3606.jpg}{
   Image of antiquarian books on a shelf from the Dr.~Philippe Garigue Collection.
   Credit: Taylor Tryburski, CC BY-SA 4.0.
}

When it comes to shelving books, bibliophiles and interior decorators rarely agree: vertical or horizontal---what’s the right way?
Ever skeptical of orthodoxy, you embrace both!
Your bookshelf displays a composed mix: upright volumes standing with quiet discipline next to a single, neatly arranged stack, like a literary ziggurat.
The effect signals bohemian chic with a touch of endearing academic absent-mindedness and just enough polish to suggest you’ve definitely read some of the books.

\medskip

A book is characterised by its spine height and thickness.
It can fit upright if its height does not exceed the bookshelf's height.
Alternatively, books can form a stack by being laid flat on their sides on top of each other; 
for aesthetic reasons the books forming such a stack must be arranged in nonincreasing order of spine height.
Their total thickness may not exceed the height of the bookshelf.
The total thickness of the upright books and the width of the stack may not exceed the width of the bookshelf.

The books in Sample~$1$ can be arranged like this:

\medskip\noindent
\centerline{\includegraphics{sample.pdf}}

